\subsection{HLD setup guide (array-based template)}

This guide assumes an array-based HLD implementation exposing \texttt{g}, \texttt{par}, \texttt{dep}, \texttt{in}, \texttt{sz}, \texttt{nxt}, \texttt{dfs\_sz}, \texttt{dfs\_hld}, \texttt{path\_query}, and a segment tree. It is separate from KACTL's class-based \texttt{HLD.h} below.

\subsubsection{Expected input}
\begin{lstlisting}
n
a[0..n-1]               // node values
n - 1 lines: u v        // undirected, 1-indexed edges
q, then q queries
\end{lstlisting}

\subsubsection{Rooting, build, and queries}
Put the following flow in \texttt{main}, after the template. \texttt{adj} stores the undirected graph; \texttt{g} stores only directed parent-to-child edges.

\begin{lstlisting}
vector<int> adj[N]; // Undirected graph; separate from g[] (children only)

cin >> n;
// Read original_value[0..n-1].
for (each edge u, v) {
    --u; --v;
    adj[u].push_back(v);
    adj[v].push_back(u);
}

// Root at 0 with BFS: convert undirected adj[] into children-only g[].
vector<int> p(n, -1);
queue<int> q;
q.push(0); p[0] = 0;
while (!q.empty()) {
    int v = q.front(); q.pop();
    for (int u : adj[v]) if (p[u] == -1) {
        p[u] = v;
        g[v].push_back(u);
        q.push(u);
    }
}

// Build HLD.
t = 0; par[0] = 0; dep[0] = 0;
dfs_sz();
nxt[0] = 0;
dfs_hld();

// Fill the segment tree using HLD positions.
for (int v = 0; v < n; v++) val[in[v]] = original_value[v];
SegTree seg(n);
seg.build();
\end{lstlisting}

First convert every one-indexed query vertex to zero indexing. The usual operations are:
\begin{lstlisting}
path query u v  -> path_query(u, v, seg)
update node x w -> seg.upd(in[x], w)
subtree of v    -> seg.qry(in[v], in[v] + sz[v] - 1)
lca(u, v)      -> lca(u, v)
distance(u, v) -> dep[u] + dep[v] - 2 * dep[lca(u, v)]
\end{lstlisting}

\subsubsection{Edge-weight version}
For input edges \texttt{u v w}, save \texttt{eu[i]}, \texttt{ev[i]}, and \texttt{ew[i]} while reading. After \texttt{dfs\_hld()}, store each edge at the position of its deeper endpoint:

\begin{lstlisting}
for (each edge i) {
    int c = dep[eu[i]] > dep[ev[i]] ? eu[i] : ev[i];
    val[in[c]] = ew[i]; // c is the child endpoint
}
// val[in[0]] remains the identity: the root has no parent edge.
\end{lstlisting}

For an edge path query, use the final segment \texttt{seg.qry(in[u] + 1, in[v])} and skip it when \texttt{u == v}. To update edge \texttt{i}, compute its child endpoint \texttt{c} in the same way and call \texttt{seg.upd(in[c], new\_w)}.

\subsubsection{Checklist}
\begin{itemize}
\item The template is zero-indexed and roots at node 0. Subtract 1 from one-indexed input.
\item For a different root \texttt{r}, start BFS at \texttt{r} and replace the root-0 initialization in \texttt{p}, \texttt{par}, \texttt{dep}, \texttt{nxt}, and the default arguments of \texttt{dfs\_sz}/\texttt{dfs\_hld}.
\item For multiple test cases, clear \texttt{g[v]} and \texttt{adj[v]} for \texttt{v} in \texttt{[0, n)} and reset \texttt{t = 0}.
\item \texttt{dfs\_sz} and \texttt{dfs\_hld} are recursive. Test a path graph with $n = 2 \cdot 10^5$ for stack overflow.
\item For max, min, xor, or gcd queries, change \texttt{merge()} and the identity value.
\item Decide whether values are on nodes or edges before setup: the two initializations differ.
\end{itemize}
